%% ma: L12-B12-0012
%% lop: 12
%% chuong: 4
%% bai: 12
%% dang: 12.12.7
%% loai: DS
%% muc: [TH, TH, VD, VD]
%% dap_an: "ĐSĐS"
%% nguon: "(NBV)-ĐỀ SỐ 10-MỖI NGÀY 1 ĐỀ THI 2025.pdf (D10, MỖI NGÀY 1 ĐỀ - Đề số 10), Phần II Câu 2"
%% kien_thuc: "Bài 12 (tích phân, quãng đường từ vận tốc); chuyển động thẳng biến đổi đều"
%% trang_thai: chinh-thuc
%% ghi_chu: "Đã duyệt 11/10/2026. Hình gốc là ảnh minh họa ô tô đâm vào khúc gỗ (trang trí, không chứa dữ kiện toán) nên không vẽ hình."
\begin{ex}
Một ô tô bắt đầu chuyển động thẳng nhanh dần đều với tốc độ $v(t)=5t$ (m/s); trong đó $t$ là thời gian tính bằng giây kể từ khi ô tô bắt đầu chuyển động. Đi được $6$ (s) người lái xe phát hiện chướng ngại vật và phanh gấp, ô tô tiếp tục chuyển động chậm dần đều với gia tốc $a=-5$ (m/s$^2$).
\choiceTF
{\True Tốc độ của ô tô tại thời điểm $10$ (s) tính từ lúc xuất phát là $10$ (m/s).}
{Quãng đường ô tô chuyển động được trong $6$ (s) đầu là $80$ m.}
{\True Quãng đường $S$ (đơn vị: mét) mà ô tô chuyển động được kể từ lúc bắt đầu đạp phanh đến khi dừng lại được tính theo công thức $S=\displaystyle\int_0^6(30-5t)\,\mathrm dt$.}
{Quãng đường ô tô chuyển động được kể từ lúc bắt đầu chuyển động cho đến khi dừng lại là $170$ m.}
\loigiai{Tại $t=6$ ô tô có tốc độ $v(6)=30$ m/s. Từ lúc phanh (đặt lại gốc thời gian), $v_1(t)=30-5t$.

a) Thời điểm $10$ s ứng với $4$ s sau khi phanh: $v_1(4)=30-20=10$ m/s. Đúng.

b) $\displaystyle\int_0^65t\,\mathrm dt=\dfrac{5t^2}2\Big|_0^6=90$ m. Sai.

c) Ô tô dừng khi $30-5t=0\Leftrightarrow t=6$, nên $S=\displaystyle\int_0^6(30-5t)\,\mathrm dt$. Đúng.

d) $S=\displaystyle\int_0^6(30-5t)\,\mathrm dt=90$ m, nên cả quãng đường là $90+90=180$ m, không phải $170$ m. Sai.}
\end{ex}
